\section{Flows, cuts and disjoint paths}\label{fl:flows}

A network has $n$ vertices, distinct terminals $s,t$, and $m$ numbered
arcs. Arc $e$ has a tail $u_e$, a head $v_e$, and a natural capacity
$c_e$. Arc numbers distinguish parallel and opposite arcs. A flow of value
$k\in\mathbb N$ assigns integers $0\le f_e\le c_e$. Its outflow equals
its inflow at each internal vertex, its outflow at $s$ equals its inflow
plus $k$, and its inflow at $t$ equals its outflow plus $k$. A cut is a
vertex mask $S$ containing $s$ and omitting $t$. Write $\delta^+(S)$ for
arcs leaving $S$ and $\delta^-(S)$ for arcs entering it. For a set of
arcs $D$, write $f(D)=\sum_{e\in D}f_e$, and similarly for $c$.

\begin{theorem}[Integral max-flow min-cut]\leavevmode\label{fl:main}
A primitive-recursive algorithm produces a flow $f$ and cut $S$ with
$k=c(\delta^+(S))$. Thus $f$ has maximum value and $S$ has minimum
capacity.
\end{theorem}
\begin{proof}
Summing the balance equations over any cut $U$ cancels every internal
arc, including loops, and gives
\[
             k+f(\delta^-(U))=f(\delta^+(U))\le c(\delta^+(U)).
\]
In particular $k\le C=c(\delta^+(\{s\}))$.
Start with the zero flow. Make a directed residual step along arc $e$
in its original direction when $f_e<c_e$, and in reverse when $f_e>0$.
Keep both the index and the direction with each step. To test whether
$t$ can be reached, enumerate simple paths from $s$ having at most
$n-1$ steps. This gives the entire reachable set: a repeated vertex in
a walk permits deletion of the intervening segment. Iterating at most
the input length leaves a simple path with the same ends. Thus the set
is closed under a further residual step.

If a simple residual $s$--$t$ path is found, add one on each forward
arc and subtract one on each reverse arc on that path. Its vertices are
distinct, so an original arc cannot receive two updates. Either type of
step changes the net outflow by $+1$ at its starting vertex and $-1$ at
its ending vertex. The changes cancel at internal vertices. The
capacity tests guarantee feasibility, and the value increases by one.
If the search fails, retain the state. Perform exactly $C$ rounds.
After the first failure, the state is fixed. If all $C$ rounds
succeeded, the value is $C$ and another success would violate the cut
bound. This also handles $C=0$. In the final residual graph let $S$ be
the reachable set. It excludes $t$. Every original arc leaving $S$
is saturated, since otherwise its forward residual step would leave
$S$. Every entering arc carries zero, since otherwise its reverse
step would leave $S$. Substitution in the displayed balance identity
gives equality. Applying its inequality to all other feasible flows
and all other cuts proves both extremal claims.
\end{proof}

A residual arc sequence of length $\ell<n$ uses digits in base $2m+2$,
with length stored separately. Flows use $m$ digits in base
$2+\max(\{0\}\cup\{c_e:e<m\})$. A cut uses $n$ bits. At round
$j\le C$, the Boolean invariant tests feasibility and the disjunction
of value $j$ with failure of the residual search. Thus the running-time
bound may depend on the numerical capacities while the induction is
quantifier-free. For nonnegative rational capacities $a_e/b_e$ with
$b_e>0$, multiply all capacities by $\prod_e b_e$, solve the integral
instance and divide the result by the product. The same cut inequality
compares the answer to every rational competitor after clearing
denominators.

\begin{lemma}[Extracting paths]\leavevmode\label{fl:decompose}
From any integral flow of value $k$ one can extract $k$ simple
$s$--$t$ paths such that each arc $e$ is used at most $f_e$ times.
\end{lemma}
\begin{proof}
For positive value, let $U$ be the vertices reachable along arcs of
positive flow. If $t\notin U$, no positive flow leaves the cut $U$,
contradicting the balance identity. Find a simple such path and
subtract one on each of its arcs. All internal balances persist and
the value decreases by one. Repeat $k$ times, retaining the list of
paths. It is an invariant that the original flow on each arc equals
the remaining flow plus the number of recorded uses. At value zero
there may still be a circulation. The stated inequality suffices.
\end{proof}

\begin{theorem}[Finite Menger theorems]\leavevmode\label{fl:menger}
For distinct $s,t$, the maximum number of pairwise arc-disjoint
directed $s$--$t$ paths equals the minimum size of a set of arcs whose
deletion destroys every such path. In an undirected graph the same
holds for edges. If the directed arc $s\to t$ is absent, the maximum
number of internally vertex-disjoint directed $s$--$t$ paths equals
the minimum size of an internal vertex separator. The undirected
vertex version holds for nonadjacent terminals.
\end{theorem}
\begin{proof}
Give each arc capacity one. The flow supplies the paths by
Lemma~\ref{fl:decompose}. The forward arcs of a minimum cut form
a separating set of the same size. Every separating set must contain
a distinct arc from each disjoint path. For an undirected edge use
two opposite unit arcs. If both carry one, subtract one from each.
This preserves the balances, so extraction uses each undirected edge
at most once. Exactly one directed arc per crossing edge leaves the
cut and those undirected edges separate.

For the vertex theorem, ignore loops, arcs entering $s$ and arcs leaving
$t$, since simple terminal paths omit them. Replace each internal
$v$ by $v^-,v^+$ and an arc $v^-\to v^+$ of capacity one. Represent each
original arc by an arc of capacity $n$ from its tail's output copy to
its head's input copy, leaving $s,t$ unsplit. Deletion of all $n-2$
unit arcs separates the terminals by the hypothesis on $s\to t$.
Take the reachable side after this deletion. Every arc leaving it in
the full network belongs to the deleted set. Hence the minimum cut
costs at most $n-2$, and none of its forward arcs has capacity $n$.
Its crossing unit arcs give a vertex separator of its capacity, since
every original path lifts through the split vertices. The extracted
paths project to simple original paths: visiting an internal vertex
uses its two copies consecutively and consumes its unit arc. Thus
their internal vertices are pairwise disjoint. Again any separator
must meet the paths at distinct vertices. For an undirected graph,
first replace each edge by two opposite arcs. All auxiliary networks
have at most $2n$ vertices and $2m+n$ arcs for $m$ original edges or arcs.
\end{proof}

\section{Matching, covering and decomposition}\label{fl:matching}

Let $(L,R,E)$ be a finite bipartite graph with disjoint parts. Put
$a=|L|$, $b=|R|$. A matching is an edge mask with at most one selected
edge incident at each vertex. A vertex cover is a vertex mask meeting
each edge. Write $N(X)\subseteq R$ for the vertices adjacent to
some $x\in X\subseteq L$. The following equality computes both an
optimal matching and a most deficient subset.

\begin{theorem}[K\H{o}nig, Hall and the deficiency formula]\leavevmode\label{fl:hall}
If $\nu$ is the maximum matching size and $\tau$ the minimum vertex-cover
size, then
\[
 \nu=\tau,\qquad
 a-\nu=\max_{X\subseteq L}\bigl(|X|-|N(X)|\bigr).
\]
Thus a matching saturates $L$ exactly when $|N(X)|\ge |X|$ for every
$X\subseteq L$. When $a=b$ such a matching is perfect. A failed
saturation test returns a specific violating $X$.
\end{theorem}
\begin{proof}
Use a source joined by unit arcs to $L$, arcs of capacity $a+1$
for each edge from $L$ to $R$, and unit arcs from $R$ to a sink.
In an integral flow the total through each internal vertex is at
most one. The positive middle arcs therefore give a matching of
size the flow value, and each matching gives such a flow. For a
minimum cut write $X$ for the left vertices on the source side and
$Y$ for the right vertices there. Its cost is at most $a$ by the
source-only cut, so no middle arc crosses and $N(X)\subseteq Y$.
Its capacity is $a-|X|+|Y|$. Replacing $Y$ by $N(X)$ preserves the
absence of crossing middle arcs. Minimality gives $|Y|=|N(X)|$.
For any other left subset $W$, the cut with sides specified by
$W,N(W)$ has cost $a-|W|+|N(W)|$. Minimizing these values gives
the formula. Interpret the differences in the formula as signed
differences, or move the neighborhood counts across the comparisons.
The empty set has deficiency zero.

The vertices $(L\setminus X)\cup Y$ meet every edge. Their number is
the minimum cut capacity, which is the maximum matching size.
Every cover contains at least one distinct vertex for each edge of
any matching. This proves $\nu=\tau$. The constructed $X$ attains
the maximum deficiency. A matching saturates the left side exactly
when its size is $a$, giving both alternatives of Hall's theorem.
The reasoning allows parallel edges, since the unit arcs at the two
ends still enforce the matching condition.
\end{proof}

\begin{corollary}[Regular bipartite graphs and edge coloring]\leavevmode\label{fl:edgecolour}
A finite $d$-regular bipartite multigraph with $d>0$ has a perfect
matching and splits into $d$ perfect matchings. Every finite bipartite
multigraph of maximum degree $D$ splits into $D$ matchings, and its
edge chromatic number is $D$ (zero for an empty edge set).
\end{corollary}
\begin{proof}
Counting edges from each side gives $da=db$, whence $a=b$.
The $d|X|$ incidences from a left subset $X$ arrive in $N(X)$,
whose total degree is $d|N(X)|$. Positive cancellation proves Hall's
condition. Removing a perfect matching lowers every degree by one.
Repeat $d$ times, testing that the remaining graph is regular of the
remaining degree and that removed edges occur exactly once in the
recorded matchings.

For the general case add isolated vertices to equalize the sides at
size $p=\max(a,b)$. If $m=|E|$, the total degree deficit to $D$ on
each side is $pD-m$. Whenever it is positive, take the first
deficient vertex on each side and add an edge between them, permitting
parallel edges with distinct indices. After $pD-m$ steps the graph is
regular. Restrict the resulting $D$ matchings to the original edges.
Conversely $D$ edges sharing an endpoint require distinct colors.
For $D=0$ the empty list of matchings suffices.
\end{proof}

\begin{theorem}[Berge's criterion and the edge-cover identity]\leavevmode\label{fl:berge}
In any finite loopless undirected graph, a matching $M$ is maximum
exactly when there is no simple path starting and ending at distinct
unmatched vertices whose edges alternate outside and inside $M$,
starting and ending outside. If the graph has $n$ vertices and no
isolated vertex, its minimum edge-cover size is $n-\nu$, where an
edge cover is an edge set incident with every vertex.
\end{theorem}
\begin{proof}
Toggling membership along the stated path matches its previously
free ends and keeps each internal vertex matched once. This increases
the size by one. Conversely take any larger matching $M'$, obtainable
by scanning the $2^m$ masks if one exists. In the symmetric difference
each vertex has degree at most two, with incident edges of opposite
types wherever the degree is two. Each nonempty edge component is an
alternating path or cycle. From a degree-one vertex follow the
unique unused continuation. If all degrees are two, start anywhere
and return to the start. The traversal uses at most $m$
edges, and no edge can attach to an internal traversed vertex after
both its incidences have been used. Cycles have equal counts of each
type. A path with more edges of $M'$ than of $M$ must occur because
the total difference is positive. Its end edges are in $M'$.
Each end is unmatched by $M$, since an incident $M$-edge would
either extend the symmetric difference there or coincide with an
$M'$-edge at an already occupied endpoint.

Take a maximum matching. No two unmatched vertices are adjacent.
Choose one incident edge for each unmatched vertex. Each chosen edge
runs to a matched vertex, and all choices are distinct. Along with
the matching they form an edge cover of size $\nu+(n-2\nu)=n-\nu$.
Conversely, in the graph formed by any edge cover of size $h$, every
component has an edge. A component on $p$ vertices has at least $p-1$
edges: start from one vertex and add an edge to a new vertex until
the component is exhausted. Summing gives at least $n-h$ components.
Choosing one edge in each produces a matching, so $\nu\ge n-h$.
Both constructions also give zero on the empty graph.
\end{proof}

\begin{corollary}[Four families of perfect graphs]\leavevmode\label{fl:perfect}
A simple bipartite graph, its complement, its line graph and the
complement of its line graph are perfect. Here perfection means that
every induced subgraph has chromatic number equal to clique number,
with both numbers zero on an empty vertex set.
\end{corollary}
\begin{proof}
Every induced subgraph of a bipartite graph remains bipartite.
If it has an edge, both numbers are two: its two parts color it
and any three vertices include two in the same part. For a nonempty
edgeless graph both numbers are one. In the complement of a bipartite
graph $H$ on $p$ vertices, a color class is a clique in $H$, hence
a singleton or an edge. The disjoint two-vertex classes form a matching.
A maximum matching therefore gives the minimum number $p-\nu(H)$ of
classes. A clique in the complement is an independent set in $H$.
Complements exchange independent sets with vertex covers, so its
maximum size is $p-\tau(H)=p-\nu(H)$. This works after any vertex
restriction.

In a line graph vertices are edges, adjacent when they meet. Inducing
on any set of edges amounts to discarding all other edges of $H$.
All pairwise incident edges of a simple bipartite graph share a
vertex if the edge family is nonempty. If two distinct edges have a common
end $v$, an edge meeting both away from $v$ would join their other ends
in the same part. Thus the line graph has clique number the maximum
degree, equal by Corollary~\ref{fl:edgecolour} to its chromatic number.
With at most one edge the assertion follows directly. In the complement
line graph a clique is a matching in $H$. A color class consists of
edges at a common vertex. The centers of the nonempty classes form a
vertex cover. Conversely assign each edge to its first endpoint in a
minimum cover to obtain such classes. Its chromatic number is therefore
$\tau(H)$, and its clique number is $\nu(H)$. Apply
Theorem~\ref{fl:hall}.
\end{proof}

\begin{theorem}[Finite Dilworth and Mirsky dualities]\leavevmode\label{fl:poset}
In a finite partially ordered set the minimum number of chains in a
partition equals the maximum size of an antichain, and the minimum
number of antichains in a partition equals the maximum size of a
chain. Extremal partitions and extremal subsets are computable.
\end{theorem}
\begin{proof}
Use two copies of the $n$ elements and join $x_L$ to $y_R$ precisely
when $x<y$. Orient each edge of a maximum matching as an arrow $x\to y$
on the original set. At most one arrow enters and at most one leaves
each element. There is no cycle by strict transitivity. Following
these arrows from all vertices with no predecessor partitions the
vertices into directed paths, including isolated vertices. Every
vertex is reached: following predecessors for $n$ steps without
finding a start would repeat a vertex. Each path is a chain, and the
number of paths is $n-\nu$ because a $k$-vertex path uses $k-1$
arrows. From a minimum vertex cover in the bipartite graph take all
elements for which neither copy belongs to the cover. At least
$n-\tau=n-\nu$ elements survive, and no two are comparable,
since their comparison edge would be uncovered. Any antichain meets
any chain at most once. Thus the displayed chain partition and this
antichain have equal sizes and are optimal.

For the dual statement enumerate all increasing sequences of length
at most $n$. For each $x$ take the largest length $r(x)$ among those
ending at $x$, including the singleton. If $x<y$, appending $y$
to a longest sequence ending at $x$ proves $r(y)>r(x)$. Equal-rank
classes are antichains. If the longest chain has length $h$, the
ranks give a partition into at most $h$ antichains, and every
antichain partition needs at least $h$ classes to cover this chain.
The empty order has both optima zero.
\end{proof}

\begin{theorem}[Stable bipartite matching]\leavevmode\label{fl:stable}
Suppose each vertex strictly ranks its acceptable neighbors, and the
edge list records mutual acceptability. Everyone prefers an acceptable
partner to being unmatched. A primitive-recursive procedure produces a
stable matching, meaning one for which no unused edge would improve both
of its endpoints. Each left vertex receives its best outcome among all
stable matchings.
\end{theorem}
\begin{proof}
Take the first free left vertex with an untried neighbor and have it
propose to its highest-ranked untried neighbor. The recipient keeps the
better of its present partner and the proposer, or accepts if free.
Mark the edge as tried and free the rejected vertex if there was one.
Held edges form a matching after every step. Pad the computation to
$m$ rounds by retaining a state with no proposal available. Each
actual proposal consumes a distinct edge, so the final state has no
free proposer with an untried neighbor. A left vertex preferring $y$
to its final outcome must already have proposed to $y$: proposals
follow preference order, and a free vertex has exhausted its list.
The recipient rejected it, either at once or later when a better
proposal arrived. A recipient's held partners only improve. Thus
$y$ prefers its final partner to that left vertex and the pair cannot
block.

Fix any stable matching $Q$. Suppose a first rejection of a pair
$(x,y)$ in $Q$ occurred. At that time $y$ keeps some $z$ preferred
to $x$. To have proposed to $y$, vertex $z$ must already have lost
every partner it ranks above $y$. None of those earlier rejections
can have been of a pair of $Q$ by the choice of the first rejection.
In $Q$, vertex $z$ is either free or paired below $y$, since $y$ is
paired with $x$. Thus $(z,y)$ blocks $Q$, a contradiction. The first
rejection is found by a scan of the finite proposal log. Therefore
no pair of $Q$ is rejected. If any left vertex preferred its partner
in $Q$ to its final outcome, it would have proposed there and been
rejected. This proves the claim simultaneously for all left vertices.
\end{proof}

\paragraph{Lean implementation detail.}
Use a digit vector for the capacities and flow, and bit masks for
cuts and matchings. A residual token is $2e$ (forward) or $2e+1$
(reverse). Its availability is a capacity comparison after division
by two. The path predicate checks the start and end vertices, that
successive token endpoints agree, that all tokens are available, and
that the $\ell+1$ visited vertices are pairwise distinct. A
length-zero path is allowed at its starting vertex. Define
\texttt{reachable}, \texttt{pathTo}, and \texttt{round} by the bounded
search and update in Theorem~\ref{fl:main}. Prove the balance identity
by induction on the arc-sum length and cancellation of the vertex
equations. For path updates, first prove the local change at each
vertex by induction through a prefix of the path, retaining the two
end corrections. Induct on the number of rounds with the stated
feasibility/value invariant. These lemmas yield equality with the
computed cut and inequalities for any competing input codes.
The remaining constructions are compositions with this solver and
bounded loops: path extraction loops on the input value, regular
completion on the common deficit, decomposition on the remaining
degree, and poset traversal on $n$. The stable-matching state contains
a tried-edge mask and a matching mask, each below $2^m$.
Each loop can check its whole finite output list by bounded digit tests.
